\begin{env}[4.3.5]
\label{0.4.3.5-fr}
Sea $\mathcal{J}$ un haz de ideales de $\mathcal{B}$. Como el funtor
$\psi^*$ es exacto, el $\psi^*(\mathcal{B})$-Módulo $\psi^*(\mathcal{J})$
se identifica canónicamente con un haz de ideales de $\psi^*(\mathcal{B})$.
La inyección canónica $\psi^*(\mathcal{J})\to\psi^*(\mathcal{B})$ da
entonces un homomorfismo de $\mathcal{A}$-Módulos
\[
  \Psi^*(\mathcal{J})
  =\psi^*(\mathcal{J})\otimes_{\psi^*(\mathcal{B})}\mathcal{A}_{[\theta]}
  \longrightarrow\mathcal{A}~;
\]
denotaremos por $\Psi^*(\mathcal{J})\mathcal{A}$, o por
$\mathcal{J}\mathcal{A}$ si no hay riesgo de confusión, la imagen de
$\Psi^*(\mathcal{J})$ por este homomorfismo. Se tiene, pues, por definición
\[
  \mathcal{J}\mathcal{A}
  =\theta^\sharp\bigl(\psi^*(\mathcal{J})\bigr)\mathcal{A}
\]
y, en particular, para todo $x\in X$,
\[
  (\mathcal{J}\mathcal{A})_x
  =\theta_x^\sharp(\mathcal{J}_{\psi(x)})\mathcal{A}_x,
\]
teniendo en cuenta la identificación canónica de los tallos de
$\psi^*(\mathcal{J})$ con los de $\mathcal{J}$ (3.7.2). Si
$\mathcal{J}_1$, $\mathcal{J}_2$ son dos ideales de $\mathcal{B}$, se tiene
\[
  (\mathcal{J}_1\mathcal{J}_2)\mathcal{A}
  =\mathcal{J}_1(\mathcal{J}_2\mathcal{A})
  =(\mathcal{J}_1\mathcal{A})(\mathcal{J}_2\mathcal{A}).
\]
Si $\mathcal{F}$ es un $\mathcal{A}$-Módulo, se pondrá
$\mathcal{J}\mathcal{F}=(\mathcal{J}\mathcal{A})\mathcal{F}$.
\end{env}
